% page 232 of the facsimile (image p-231 of the scan)
% block 165.1 x 242.0 mm ; left margin 8.0 mm
\pgc{-12.700mm}{0.000mm}{
\l{\cel{5}224.}
\l{}
\l{\sou{hostio}. (\sou{liturgio katolika}). Pano-peco dina, sen-hefa, desti-}
\l{\cel{3}nita a la meso-sakrifiko. (La\cel{1}\sou{santa hostio}\cel{1}: la pano konsa-}
\l{\cel{3}krita e chanjita aden la korpo di Iesu Kristo). - EF.}
\l{}
\l{\sou{hotelo}.\cel{2}Domo diqua onu lugas la chambri,la apartmenti, mobloza}
\l{\cel{3}ed ube onu restoras la voyajeri. - DEFR.}
\l{}
\l{\sou{\textquotedbl{}hozanna\textquotedbl{}}. I. (\sou{liturgio Hebrea}). Refreno di himno di la sinagogo.}
\l{\cel{3}- II. (\sou{liturgio katolika)}. Himno quan onu kantas en la Pal-}
\l{\cel{3}mo-sundio,qua komencas per la vorto \textquotedbl{}Hozanna\textquotedbl{}.}
\l{}
\l{hu !\cel{2}Interjeciono uzata por shamigar, por mokar, ulu.}
\l{}
\l{\sou{huo}. (\sou{zool.})\cel{2}Rapt-ucelo noktala quan onu dresis por chasado.}
\l{\cel{3}- L. athene noctua.\cel{2}- F.}
\l{}
\l{\sou{hufo}. Envelopilo korno-harda qua cirkondas la pedo di la rumi-}
\l{\cel{3}neri, pakidermi, solipedi. - DE.}
\l{}
\l{hugenoto. (Uzita desprizoze da la katoliki).\cel{2}Partisano di la}
\l{\cel{3}herezio da Calvin.\cel{2}- DEFIRS.}
\l{}
\l{\sou{hum}\cel{1}!\cel{2}Interjeciono por indikar desfido, dubito.}
\l{}
\l{\cel{1}\sou{humanismo}. Doktrino di la +humanisti di Renesanco, qui igis}
\l{\cel{3}tre prizata itere la lingui e la literaturo-verki di la}
\l{\cel{3}Grekia e la Roma antiqua. - DEFIRS.}
\l{}
\l{\cel{1}\sou{humanisto}. Persono qua par-savas la lingui,ed esas erudita}
\l{\cel{3}pri la literaturo-verki, di la Grekia e la Roma antiqua.}
\l{\cel{3}- DEFIRS.}
\l{}
\l{\sou{humero}. (\sou{anat.}) Osto di la brakio, de la shultro til la kudo.}
\l{\cel{3}- EFIS.}
\l{}
\l{\sou{humida.}\cel{2}Qua kontenas en sua materio substanco aqua od aquatra.}
\l{\cel{3}- EFIS.}
\l{}
\l{\sou{humila}.- Qua su abasas volunte koram la kunhomi. - EFIS.}
\l{}
\l{\sou{humoro}. I. Liquido qua renkontresas en korpo organika. -}
\l{\cel{3}II. Ta humoro egardata kom alterita e kom la kauzo di ula}
\l{\cel{3}morbi. - III. Temperamento, karaktero. - EFIRS.}
\l{}
\l{\sou{humuro}. Gayeso flegmoza, ulakaze akompanata da ironio.-DEFIRS.}
\l{}
\l{\sou{humuso.}\cel{2}Materio bruna o nigratra, qua produktesas en sulo,}
\l{\cel{3}da la deskompozeso di la materii organika. - DEFIS.}
\l{}
\l{\sou{hundo}. (\sou{zool.}) Genero de mamifero digitigrada, domestika, qua}
\l{\cel{3}aboyas ; e di qua existas multa varietati. - L.\cel{1}\sou{canis}}
\l{\cel{3}familiaris. - DE.}
\l{}
\l{\sou{hungrar}. (\sou{trans., ulo)}. Bezonar manjo. - DF.}
}
